\section{Operator-ideal details}
\label{app:operator-details}

\subsection{A compactness lemma tailored to symmetric rows}

Let $A=(a_{mn})$ be a symmetric coefficient matrix and suppose its absolute
row sums $r_m=\sum_n|a_{mn}|$ are finite, uniformly bounded, and tend to
zero.  The symmetric Schur test gives $\|A\|\leq\sup_m r_m$.  Moreover,
\[
 \|Q_MAQ_M\|\leq\sup_{m>M}r_m\longrightarrow0.
\]
The operator $A-Q_MAQ_M=P_MA+Q_MAP_M$ is finite rank: the first term has
finite-dimensional range and the second has finite-dimensional initial
space.  Hence $A$ is compact.  This argument is the precise ``finite tail
control'' used after \eqref{eq:row-decay}; it does not assume compactness in
advance.

For a compact operator $K$ and projections $P_N\to I$ strongly,
$(I-P_N)K\to0$ and $K(I-P_N)\to0$ in norm.  Therefore
\[
 \|K-P_NKP_N\|
 \leq\|(I-P_N)K\|+\|P_NK(I-P_N)\|\longrightarrow0,
\]
which proves \eqref{eq:compression-convergence} in the exact compression
form quoted in the main text.

\subsection{Phase removal and its limit}

For $s=\sigma+i\tau$, define
$U_\tau=\diag(n^{-i\tau/2})_{n\geq1}$.  Since $|n^{-i\tau/2}|=1$,
$U_\tau$ is unitary, and direct multiplication of matrix entries gives
\[
 (U_\tau E_\sigma U_\tau)_{mn}
 =m^{-i\tau/2}e_\sigma(m,n)n^{-i\tau/2}=e_s(m,n).
\]
If $E_\sigma$ is bounded, this matrix identity on finitely supported
vectors extends by continuity.  Conversely,
$E_\sigma=U_\tau^*E_sU_\tau^*$, so boundedness and every unitarily invariant
singular-value ideal are equivalent for $s$ and $\sigma$.

The two occurrences of $U_\tau$ point in the same direction.  In general
$U_\tau\neq U_\tau^*$, and consequently this calculation says nothing
about eigenvalues.  The direct complex trace calculations in
\cref{sec:traces-determinants} are therefore necessary.

\subsection{Absolute entry summability is nuclear}

For clarity, the trace-class sufficiency argument uses the expansion
\[
 E_s=\sum_{m,n}e_s(m,n)(e_m\otimes e_n^*).
\]
Each matrix unit has trace norm one.  If
$\sum_{m,n}|e_s(m,n)|<\infty$, the series converges in trace norm and
$\|E_s\|_1$ is bounded by that sum.  Equation
\eqref{eq:absolute-entry-bound} proves exactly this hypothesis for
$\sigma>1$.  Necessity comes from the diagonal inequality and is not the
converse of entrywise summability.

Similarly, the Hilbert--Schmidt norm is exactly the $\ell^2$ norm of all
matrix entries.  The loop subseries proves necessity at $1/2$, whereas the
coprime double majorant proves sufficiency.  No interpolation argument is
needed, and no assertion is made for all $\Schatten p$.

